# Cover Letter

Dear Editor,

We submit our manuscript entitled "When is the q-Weibull tail parameter identifiable under interval censoring? A cautionary tale with frailty equivalence" for consideration in *Communications in Statistics - Simulation and Computation*.

The q-Weibull distribution is increasingly used in reliability and survival analysis, and researchers routinely interpret an estimated $\hat q>1$ as evidence of a heavier-than-Weibull tail. We show that this interpretation is not identifiable under interval censoring: the q-Weibull survival function is mathematically close to the Weibull with gamma frailty, and the two models achieve identical log-likelihoods on three real interval-censored datasets (dental emergence, $n=4386$; diabetes onset, $n=731$; mouse tumour latency, $n=144$).

Our contributions are:
1. A Fisher information analysis showing that $q$ has only 7% independent information after profiling out scale and shape.
2. A 200-replication simulation demonstrating bidirectional confounding between q-Weibull and frailty models.
3. Three real-data applications confirming the equivalence when $q>1$.
4. Practical diagnostics for when $q$ is identifiable ($q<1$ bounded support, clustered data, large samples).

This is a simulation-focused paper that we believe fits the scope of *Communications in Statistics - Simulation and Computation* well. The code and datasets are available at https://github.com/fzmath/qweibull-ic.

Sincerely,
Huirong Cao
Corresponding author: huirongcao@126.com
